% !PRINTED 258
\weberpage{258}
Wir nehmen nun an, die Form $\varphi(x_1, x_2 \ldots x_m)$ lasse sich
durch $k$ von einander unabhängige lineare Formen von $x$ aus\-drücken,
die wir mit
\weq{2}{
\begin{aligned}
z_1 &= \beta_1^{(1)}x_1 + \beta_2^{(1)}x_2 + \cdots + \beta_k^{(1)}x_k + \beta_{k+1}^{(1)}x_{k+1} + \cdots + \beta_m^{(1)}x_m \\
&\cdot\ \cdot\ \cdot\ \cdot\ \cdot\ \cdot\ \cdot\ \cdot\ \cdot\ \cdot\ \cdot\ \cdot\ \cdot\ \cdot\ \cdot\ \cdot\ \cdot\ \cdot \\
z_k &= \beta_1^{(k)}x_1 + \beta_2^{(k)}x_2 + \cdots + \beta_k^{(k)}x_k + \beta_{k+1}^{(k)}x_{k+1} + \cdots + \beta_m^{(k)}x_m
\end{aligned}
}
bezeichnen wollen.

Da $z_1, z_2 \ldots z_k$ von einander linear unabhängig sein sollen, so
muss unter den aus der Matrix
\[
\begin{matrix}
\beta_1^{(1)}, & \beta_2^{(1)} & \ldots & \beta_m^{(1)} \\
\cdot & \cdot & \cdot & \cdot \\
\beta_1^{(k)}, & \beta_2^{(k)} & \ldots & \beta_m^{(k)}
\end{matrix}
\]
zu bildenden $k$-reihigen Determinanten wenigstens eine von Null
verschieden sein. Denn wären sie alle gleich Null, so liessen
sich die Unbekannten $h_1, h_2 \ldots h_k$, die nicht alle verschwinden
sollen, aus den Gleichungen
\[
h_1 \beta_s^{(1)} + h_2 \beta_s^{(2)} + \cdots + h_k \beta_s^{(k)} = 0, \quad s = 1, 2, 3 \ldots m
\]
bestimmen (\S.\,24), und die $z_1, z_2 \ldots z_k$ würden einer Gleichung
\[
h_1 z_1 + h_2 z_2 + \cdots + h_k z_k = 0
\]
genügen, also nicht unabhängig sein.

Wir können aber, ohne die Allgemeinheit zu beschränken, an\-nehmen,
dass die Determinante
\[
\Sigma \pm \beta_1^{(1)}\beta_2^{(2)} \ldots \beta_k^{(k)}
\]
von Null verschieden sei, und dann können wir die Gleichungen
(2) in Bezug auf die Variablen $x_1, x_2 \ldots x_k$ auflösen.

Diesen Auflösungen geben wir die Form
\weq{3}{
\begin{aligned}
x_1 &= y_1 + \alpha_{k+1}^{(1)}x_{k+1} + \cdots + \alpha_m^{(1)}x_m \\
x_2 &= y_2 + \alpha_{k+1}^{(2)}x_{k+1} + \cdots + \alpha_m^{(2)}x_m \\
&\cdot\ \cdot\ \cdot\ \cdot\ \cdot\ \cdot\ \cdot\ \cdot\ \cdot\ \cdot\ \cdot\ \cdot \\
x_k &= y_k + \alpha_{k+1}^{(k)}x_{k+1} + \cdots + \alpha_m^{(k)}x_m,
\end{aligned}
}
worin die $y_s$ lineare homogene Verbindungen der $z_s$ sind, die
aus den Gleichungen
\[
\begin{aligned}
\beta_1^{(1)}y_1 + \beta_2^{(1)}y_2 + \cdots + \beta_k^{(1)}y_k &= z_1, \\
\beta_1^{(2)}y_1 + \beta_2^{(2)}y_2 + \cdots + \beta_k^{(2)}y_k &= z_2, \\
&\cdot\ \cdot\ \cdot\ \cdot\ \cdot\ \cdot\ \cdot\ \cdot\ \cdot\ \cdot\ \cdot\ \cdot \\
\beta_1^{(k)}y_1 + \beta_2^{(k)}y_2 + \cdots + \beta_k^{(k)}y_k &= z_k
\end{aligned}
\]
%CONT>>
